% from ega0-18-fr.tex
\oldpage[0\textsubscript{IV-1}]{64}

où $E'=F\oplus_{(\mathfrak J/\mathfrak J^2)}L$ (18.2.8)~; il s'agit de
prouver que $u'$ est une $B$-\emph{équivalence}, ce qui établira que $\eta_L$
est bijectif. En vertu de la construction faite dans (18.2.8), tout revient à
prouver (compte tenu de (18.1.7)) que l'image réciproque $G$ de l'extension
$E$ de $C$ par l'homomorphisme canonique $g:F\to C$ est $F$-triviale, ce qui
est évident puisque $g$ se factorise en
$F\xrightarrow{u}E\xrightarrow{p}C$ (18.1.6).

Il reste à voir que $\eta_L$ est un homomorphisme de \emph{groupes}~; or, pour
tout $B$-homomorphisme $h:L\to L'$, il est immédiat que le diagramme
\[
  \xymatrix@C=3pc@R=2.4pc{
    \operatorname{Hom}_C(\mathfrak J/\mathfrak J^2,L)
      \ar[r]^{\eta_L}\ar[d]_{\operatorname{Hom}(1,h)}
      & \operatorname{Exan}_B(C,L)\ar[d]^{h_*}\\
    \operatorname{Hom}_C(\mathfrak J/\mathfrak J^2,L')
      \ar[r]_{\eta_{L'}}^{\sim}
      & \operatorname{Exan}_B(C,L')
  }
\]
est commutatif. Il suffit d'appliquer cette remarque à l'homomorphisme
$L\times L\to L$ définissant l'addition pour conclure.
\end{proof}

\subsection{Extensions d'algèbres}
\label{subsection:0.18.4-fr}

\begin{env}[18.4.1]
\label{0.18.4.1-fr}
Soit $A$ un anneau \emph{commutatif}. La catégorie des $A$-\emph{algèbres} est
alors une sous-catégorie pleine de celle des $A$-anneaux, caractérisée par le
fait que les homomorphismes structuraux $\rho:A\to B$ sont \emph{centraux}.

Si $B$ est une $A$-algèbre, $L$ un $B$-bimodule, il est immédiat que la
$A$-extension triviale type $D_B(L)$ est une $A$-algèbre. De même, dans la
construction de (18.2.5) (resp. de (18.2.8)), si $B$, $B'$ et $E'$ (resp.
$B$ et $E$) sont des $A$-algèbres, il en est de même de
$E'\times_{B'}B$ (resp. de $E\oplus_LL'$). Enfin, il est clair que si $B$ est
une $A$-algèbre et $E$ une $A$-extension de $B$ par un $B$-bimodule $L$ qui
est une $A$-algèbre, toute $A$-extension de $B$ par $L$, $A$-\emph{équivalente}
à $E$, est aussi une $A$-algèbre. On déduit alors aussitôt de la définition
(18.3.4) que les classes de $A$-extensions équivalentes de $B$ par $L$ qui
sont des $A$-\emph{algèbres} forment un \emph{sous-groupe}, noté
$\operatorname{Exal}_A(B,L)$, de $\operatorname{Exan}_A(B,L)$. Soit
$\mathbf K'$ la sous-catégorie pleine de la catégorie $\mathbf K$ définie dans
(18.3.5), dont les objets $(A,B,L)$ sont tels que $A$ soit commutatif et $B$
une $A$-algèbre. Alors ce qui précède montre que
\[
  (A,B,L)\longmapsto\operatorname{Exal}_A(B,L)
\]
est un \emph{foncteur covariant} de $\mathbf K'$ dans $\mathbf{Ab}$. Les
résultats de (18.3.7) sont inchangés lorsqu'on remplace partout
$\operatorname{Exan}$ par $\operatorname{Exal}$.
\end{env}

\begin{env}[18.4.2]
\label{0.18.4.2-fr}
Supposons toujours l'anneau $A$ commutatif. Les remarques de (18.4.1) restent
valables lorsque l'on remplace «~algèbre~» par «~algèbre commutative~» et
«~bimodule~» par «~module~». Si $B$ est une $A$-algèbre commutative et $L$ un
$B$-module,

